% ============================================================
%  CHAPTER 1 — GRAMMAR
% ============================================================
\section{Grammar}

% ------------------------------------------------------------
\subsection{1.1 Script \& Sounds}

\subsubsection{Vowels (Swar — \hw{स्वर})}

\begin{tabular}{@{}llll@{}}
  \toprule
  Roman & Devanagari & Note \\
  \midrule
  a   & \hw{अ} & /ə/ & short, like \textit{u} in \textit{but} \\
  aa  & \hw{आ} & /aː/ & long, like \textit{a} in \textit{father} \\
  i   & \hw{इ} & /ɪ/ & short, like \textit{i} in \textit{bit} \\
  ii  & \hw{ई} & /iː/ & long, like \textit{ee} in \textit{see} \\
  u   & \hw{उ} & /ʊ/ & short, like \textit{u} in \textit{put} \\
  uu  & \hw{ऊ} & /uː/ & long, like \textit{oo} in \textit{food} \\
  e   & \hw{ए} & /eː/ & like \textit{a} in \textit{say} \\
  ai  & \hw{ऐ} & /ɛː/ & like \textit{e} in \textit{bed} \\
  o   & \hw{ओ} & /oː/ & like \textit{o} in \textit{go} \\
  au  & \hw{औ} & /ɔː/ & like \textit{aw} in \textit{law} \\
  % --- add more as covered ---
  \bottomrule
\end{tabular}

\subsubsection{Consonants (Vyanjan — \hw{व्यंजन})}
\textit{[Placeholder — expand row by row as covered with teacher.]}

% ------------------------------------------------------------
\subsection{1.2 Nouns \& Gender}

Hindi has two grammatical genders: \textbf{masculine} and \textbf{feminine}. Gender affects adjective and verb agreement throughout the sentence.

\begin{itemize}
  \item Most masculine nouns end in \textbf{-aa}: \textit{ladkaa} (\hw{लड़का}) — boy
  \item Most feminine nouns end in \textbf{-ii}: \textit{ladkii} (\hw{लड़की}) — girl
  \item Many exceptions exist — always learn gender with the noun.
\end{itemize}

\subsubsection{Plurals}

\begin{tabular}{@{}llll@{}}
  \toprule
  Gender & Singular ending & Plural ending & Example \\
  \midrule
  Masculine & -aa & -e & \textit{ladkaa} → \textit{la.Rke} \\
  Feminine  & -ii & -iyaa.N & \textit{ladkii} → \textit{la.Rkiyaa.N} \\
  Other     & — & unchanged & \textit{kitaab} → \textit{kitaaben} \\
  \bottomrule
\end{tabular}

% ------------------------------------------------------------
\subsection{1.3 Pronouns (\hw{सर्वनाम})}

\begin{tabular}{@{}lllp{4cm}@{}}
  \toprule
  English & Transliteration & Hindi & Note \\
  \midrule
  I              & main    & \hw{मैं}   & \\
  you (informal) & tuu     & \hw{तू}    & intimate; use with care \\
  you (familiar) & tum     & \hw{तुम}   & friends, family \\
  you (formal)   & aap     & \hw{आप}   & default polite form \\
  he / she / it (near) & yah & \hw{यह} & \\
  he / she / it (far)  & vah & \hw{वह} & often pronounced \textit{vo} \\
  we             & ham     & \hw{हम}   & \\
  they           & ve      & \hw{वे}   & \\
  \bottomrule
\end{tabular}

\subsubsection{Oblique Pronouns (used before postpositions)}

\begin{tabular}{@{}llll@{}}
  \toprule
  Subject form & Oblique form & Hindi & Example use \\
  \midrule
  main & mujh & \hw{मुझ} & \textit{mujhe} (to me), \textit{mujhse} (from me) \\
  tuu  & tujh & \hw{तुझ} & \\
  tum  & tum  & \hw{तुम} & \textit{tumhe}, \textit{tumse} \\
  aap  & aap  & \hw{आप}  & \textit{aapko}, \textit{aapse} \\
  vah  & us   & \hw{उस}  & \textit{usse}, \textit{usko} \\
  ham  & ham  & \hw{हम}  & \textit{hamko}, \textit{hamse} \\
  ve   & un   & \hw{उन}  & \textit{unhe}, \textit{unse} \\
  \bottomrule
\end{tabular}

% ------------------------------------------------------------
\subsection{1.4 Postpositions (\hw{परसर्ग})}

Hindi uses \textbf{postpositions} placed \emph{after} the noun (in its oblique form).

\begin{tabular}{@{}llp{5.5cm}@{}}
  \toprule
  Postposition & Hindi & Meaning \\
  \midrule
  men       & \hw{में}         & in, inside \\
  par       & \hw{पर}          & on, at \\
  se        & \hw{से}          & from, with, by (instrument) \\
  ko        & \hw{को}          & to, for (indirect object / dative) \\
  kaa/kii/ke & \hw{का/की/के}  & of, 's (possession; agrees with possessed noun) \\
  tak       & \hw{तक}          & until, up to \\
  ke liye   & \hw{के लिए}     & for (purpose) \\
  ke paas   & \hw{के पास}     & near, with (possession) \\
  ke saath  & \hw{के साथ}     & with (accompaniment) \\
  ke baad   & \hw{के बाद}     & after \\
  se pehle  & \hw{से पहले}    & before \\
  % --- add more as learned ---
  \bottomrule
\end{tabular}

% ------------------------------------------------------------
\subsection{1.5 Adjectives (\hw{विशेषण})}

Adjectives ending in \textbf{-aa} agree with the noun in gender and number; all others are invariable.

\begin{tabular}{@{}lllll@{}}
  \toprule
  English & Masc. sg. & Masc. pl./Obl. & Fem. & Hindi \\
  \midrule
  good   & acchaa  & acche  & acchii  & \hw{अच्छा / अच्छे / अच्छी} \\
  big    & badaa  & bade  & badii  & \hw{बड़ा / बड़े / बड़ी} \\
  new    & nayaa   & naye   & naii    & \hw{नया / नए / नई} \\
  hot    & garam   & garam  & garam   & \hw{गरम} — invariable \\
  beautiful & sundar & sundar & sundar & \hw{सुंदर} — invariable \\
  % --- add more ---
  \bottomrule
\end{tabular}

% ------------------------------------------------------------
\subsection{1.6 Verbs (\hw{क्रिया})}

\subsubsection{Infinitive}
All Hindi infinitives end in \textbf{-naa}: e.g.\ \textit{jaanaa} (\hw{जाना}) — \textit{to go}.
Root = infinitive minus \textit{-naa}: \textit{jaa-}.

\subsubsection{The verb \textit{honaa} (\hw{होना}) — to be}

\begin{tabular}{@{}llll@{}}
  \toprule
  Subject & Present & Transliteration & Hindi \\
  \midrule
  main        & hoon & hoon  & \hw{हूँ} \\
  tuu         & hai   & hai   & \hw{है} \\
  tum         & ho    & ho    & \hw{हो} \\
  aap/vah/yah & hai   & hai   & \hw{है} \\
  ham/ve      & hain  & hain  & \hw{हैं} \\
  \bottomrule
\end{tabular}

\subsubsection{Present Habitual Tense}

\pattern{Subject + [Object] + verb-stem + taa/tii/te + hoon/hai/hain/ho}

\textit{-taa} → masculine singular \quad \textit{-tii} → feminine \quad \textit{-te} → masculine plural / formal

\begin{tabular}{@{}lll@{}}
  \toprule
  Transliteration & Hindi & English \\
  \midrule
  main chaay piitaa hoon & \hw{मैं चाय पीता हूँ.} & I drink tea. (m.) \\
  vah roz jaatii hai      & \hw{वह रोज़ जाती है.}  & She goes every day. \\
  ham Hindi sikhte hain   & \hw{हम हिंदी सीखते हैं.} & We learn Hindi. \\
  \bottomrule
\end{tabular}

\subsubsection{Present Continuous Tense}

\pattern{Subject + verb-stem + rahaa/rahii/rahe + hoon/hai/hain/ho}

\textit{-rahaa} → masculine singular \quad \textit{-rahii} → feminine \quad \textit{-rahe} → plural / formal

\begin{tabular}{@{}lll@{}}
  \toprule
  Transliteration & Hindi & English \\
  \midrule
  main Hindi sikh rahaa hoon     & \hw{मैं हिंदी सीख रहा हूँ।}       & I am learning Hindi. (m.) \\
  merii patnii khaa rahii hai    & \hw{मेरी पत्नी खाना खा रही है।}   & My wife is eating. \\
  aap kyaa kar rahe ho?          & \hw{आप क्या कर रहे हो?}            & What are you doing? \\
  bacchaa ananaas khaa rahaa hai & \hw{बच्चा अनानास खा रहा है।}       & The child is eating pineapple. \\
  \bottomrule
\end{tabular}

\subsubsection{Simple Past Tense (भूतकाल)}

\pattern{Intransitive: Subject + verb-stem + aa/ii/e + (thaa/thii/the)}
\pattern{Transitive:\ \ \ \ Subject\,+\,\textbf{ne} + Object + verb-pp + thaa/thii/the}

Transitive verbs (those taking a direct object) require \textbf{-ne} on the subject.\\
Intransitive verbs (motion, states) do \emph{not} take \textit{-ne}.

\begin{tabular}{@{}lll@{}}
  \toprule
  Transliteration & Hindi & English \\
  \midrule
  mainne painkek khaayaa thaa  & \hw{मैंने पैनकेक खाया था।}  & I ate pancakes. (transitive + ne) \\
  mainne painkek banaayaa thaa & \hw{मैंने पैनकेक बनाया था।} & I made pancakes. (transitive + ne) \\
  main aafis gayaa thaa        & \hw{मैं ऑफिस गया था।}       & I went to the office. (intransitive) \\
  vah ghar aaii                & \hw{वह घर आई।}              & She came home. (intransitive) \\
  \bottomrule
\end{tabular}

\subsubsection{Past Habitual — \textit{used to}}

\pattern{Subject + verb-stem + taa/tii/te + thaa/thii/the}

Used for repeated or habitual actions in the past. Compare: Present Habitual uses \textit{hoon/hai/hain}.

\begin{tabular}{@{}lll@{}}
  \toprule
  Transliteration & Hindi & English \\
  \midrule
  vah ek gaanv men rahtii thii      & \hw{वह एक गाँव में रहती थी।}       & She used to live in a village. \\
  vahaan bandar the                  & \hw{वहाँ बंदर थे।}                  & There were monkeys there. \\
  log vahaan shaanti se baiThte the  & \hw{लोग वहाँ शांति से बैठते थे।}   & People used to sit there quietly. \\
  \bottomrule
\end{tabular}

\subsubsection{Future Tense (भविष्य काल)}

\pattern{Subject + [Object] + verb-stem + uungaa/uungii/enge (inflects for gender \& person)}

\begin{tabular}{@{}llll@{}}
  \toprule
  Subject & Masc.\ ending & Fem.\ ending & Example \\
  \midrule
  main    & -uungaa  & -uungii  & \hw{मैं जाऊंगा / जाऊंगी।} \\
  tum     & -oge     & -ogii    & \hw{तुम जाओगे / जाओगी।} \\
  aap/vah & -enge    & -engii   & \hw{आप/वह जाएंगे / जाएंगी।} \\
  ham     & -enge    & -engii   & \hw{हम जाएंगे।} \\
  \bottomrule
\end{tabular}

\begin{tabular}{@{}lll@{}}
  \toprule
  Transliteration & Hindi & English \\
  \midrule
  main kal aafis jaaungaa  & \hw{मैं कल ऑफिस जाऊंगा।}   & I will go to the office tomorrow. \\
  kyaa aap paanii pienge?  & \hw{क्या आप पानी पिएंगे?}   & Will you drink water? \\
  vah khaanaa banaaegii    & \hw{वह खाना बनाएगी।}         & She will make food. \\
  \bottomrule
\end{tabular}

\subsubsection{Imperative}
\textit{[Placeholder — fill in when covered.]}

% ------------------------------------------------------------
\subsection{1.7 Question Words (\hw{प्रश्नवाचक शब्द})}

\begin{tabular}{@{}llll@{}}
  \toprule
  English & Transliteration & Hindi & Example \\
  \midrule
  what      & kyaa    & \hw{क्या}   & \hw{यह क्या है?} — What is this? \\
  who       & kaun    & \hw{कौन}    & \hw{वह कौन है?} — Who is that? \\
  where     & kahaan & \hw{कहाँ}   & \\
  when      & kab     & \hw{कब}     & \\
  why       & kyõ    & \hw{क्यों}  & \\
  how       & kaise   & \hw{कैसे}   & \hw{आप कैसे हैं?} — How are you? \\
  how much  & kitnaa  & \hw{कितना}  & inflects for gender \\
  \bottomrule
\end{tabular}

% ------------------------------------------------------------
\subsection{1.8 Sentence Order}

Hindi is \textbf{Subject–Object–Verb (SOV)}.

\pattern{Subject \ \ [Indirect Object + ko] \ \ [Direct Object] \ \ [Adverb] \ \ Verb}

\medskip
Compare with English (SVO):

\begin{tabular}{@{}lll@{}}
  \toprule
  Language & Order & Example \\
  \midrule
  English & S–V–O & I drink tea. \\
  Hindi   & S–O–V & \hw{मैं चाय पीता हूँ.} \\
  \bottomrule
\end{tabular}

\subsection{1.9 Obligation — \textit{ko + Infinitive + honaa}}

Used to express necessity ("must", "have to"). The subject takes oblique form + \textit{ko}; the infinitive stays unchanged.

\pattern{Subject (oblique + ko) + Infinitive + hogaa / hai}

\begin{tabular}{@{}lll@{}}
  \toprule
  Transliteration & Hindi & English \\
  \midrule
  mujhe jaanaa hogaa       & \hw{मुझे जाना होगा।}         & I will have to go. \\
  aapko Nepaal aanaa hogaa & \hw{आपको नेपाल आना होगा।}    & You must come to Nepal. \\
  mujhe 5 km dauRnaa hogaa & \hw{मुझे 5 km दौड़ना होगा।}  & I will have to run 5 km. \\
  tumhen kaam karnaa hai   & \hw{तुम्हें काम करना है।}     & You have to work. \\
  \bottomrule
\end{tabular}

\subsection{1.10 Comparison (\textit{se} + adjective)}

The reference point (what you compare \emph{against}) takes \textit{se}. The adjective agrees with the subject.

\pattern{X + Y-se + baRaa/chhoTaa (inflected) + hai}

\begin{tabular}{@{}lll@{}}
  \toprule
  Transliteration & Hindi & English \\
  \midrule
  merii bahan mujhse baRii hai    & \hw{मेरी बहन मुझसे बड़ी है।}   & My sister is bigger/older than me. \\
  meraa bhaaii mujhse chhoTaa hai & \hw{मेरा भाई मुझसे छोटा है।}  & My brother is smaller/younger than me. \\
  Roos ek baRaa desh hai          & \hw{रूस एक बड़ा देश है।}       & Russia is a big country. \\
  \bottomrule
\end{tabular}

\subsection{1.11 [Add New Grammar Topic Here]}

\textit{[Add heading, brief explanation, and example table.]}
